% GENERATED FILE. Edit canonical lessons, then run npm run generate:books.
% canonical-source-hash: fnv1a64:d51427381e72c90c
% canonical-lessons: FA-C462-haraj, FA-C462-arzidan, FA-C462-kharid-kardan, FA-C462-forush-raftan, FA-C462-jabe

\chapter{Sales}
\label{ch:fa-sales}
\hlchaptermodality{462}

\begin{chapteropening}
\textbf{By the end of this chapter:} \emph{I can talk about sales and spending: a sale, being worth it, doing the shopping, selling well and a box.}

Complete the last lesson of chapter 462: I can talk about sales and spending: a sale, being worth it, doing the shopping, selling well and a box.
\end{chapteropening}

\section[\texorpdfstring{\fa{حراج} (harâj)}{حراج (harâj)}]{\fa{حراج} (harâj) — a sale (with reduced prices); an auction}

\label{lesson:FA-C462-haraj}



\begin{quote}
Before the new one: say the Persian for small change, then the Persian for a bureau de change.
\end{quote}

\subsection*{You'll want to know: \fa{حراج}}
\textbf{\fa{حراج}} (\emph{harâj}) — \textquotedblleft{}a sale (with reduced prices); an auction\textquotedblright{}.

\begin{quote}
\fa{در} \fa{حراج} \fa{زمستانی} \fa{یک} \fa{کاپشن} \fa{ارزان} \fa{خریدم}.
\end{quote}

\begin{quote}
\emph{Dar harâj-e zemestâni yek kâpshan-e arzân kharidam.}
\end{quote}

\begin{quote}
I bought a cheap jacket in the winter sale.
\end{quote}

\emph{Harâj!} written in a shop window means a sale is on.

\begin{cousinweb}
\emph{harâj} was borrowed from Arabic and is usually traced to Arabic \emph{ḥarāj} 'an auction'. EtymDB does not list it.
\end{cousinweb}

\begin{grammarlens}[title={What you've built}]
One more word for telling, explaining and describing at B1.
\end{grammarlens}

\subsection*{Your turn}
\begin{itemize}
\raggedright
  \item \emph{Say it:} \emph{harâj}
  \item \emph{Say it:} \emph{harâj}, once more
  \item \emph{Say it:} a sentence of your own with \emph{harâj}
  \item \emph{Recall:} say the Persian for small change, then the Persian for a bureau de change, then say \emph{harâj} again
\end{itemize}

\begin{tcolorbox}[breakable,colback=teal!4,colframe=teal!35!black,title={Before you move on}]
What does \fa{حراج} mean? (\textquotedblleft{}A sale (with reduced prices); an auction\textquotedblright{}.) Say it once more.
\end{tcolorbox}
\section[\texorpdfstring{\fa{ارزیدن} (arzidan)}{ارزیدن (arzidan)}]{\fa{ارزیدن} (arzidan) — to be worth}

\label{lesson:FA-C462-arzidan}



\begin{quote}
Before the new one: say the Persian for a bureau de change, then the Persian for a sale.
\end{quote}

\subsection*{You'll want to know: \fa{ارزیدن}}
\textbf{\fa{ارزیدن}} (\emph{arzidan}) — \textquotedblleft{}to be worth\textquotedblright{}.

\begin{quote}
\fa{این} \fa{کفش} \fa{به} \fa{این} \fa{قیمت} \fa{نمیارزد}.
\end{quote}

\begin{quote}
\emph{In kafsh be in qeymat nemiarzad.}
\end{quote}

\begin{quote}
These shoes aren't worth this price.
\end{quote}

Present stem \emph{arz-}: \emph{miarzad}. Built on the same root as \emph{arzân} 'cheap'.

\begin{cousinweb}
\emph{arzidan} is inherited: EtymDB traces it through Proto-Iranian \emph{Hárǰati} to the PIE root \emph{h\textsubscript{2}elg\textsuperscript{w}\textsuperscript{h}-} 'to be worth, value'.
\end{cousinweb}

\begin{grammarlens}[title={What you've built}]
One more word for telling, explaining and describing at B1.
\end{grammarlens}

\subsection*{Your turn}
\begin{itemize}
\raggedright
  \item \emph{Say it:} \emph{arzidan}
  \item \emph{Say it:} \emph{arzidan}, once more
  \item \emph{Say it:} a sentence of your own with \emph{arzidan}
  \item \emph{Recall:} say the Persian for a bureau de change, then the Persian for a sale, then say \emph{arzidan} again
\end{itemize}

\begin{tcolorbox}[breakable,colback=teal!4,colframe=teal!35!black,title={Before you move on}]
What does \fa{ارزیدن} mean? (\textquotedblleft{}To be worth\textquotedblright{}.) Say it once more.
\end{tcolorbox}
\section[\texorpdfstring{\fa{خرید} \fa{کردن} (kharid kardan)}{خرید کردن (kharid kardan)}]{\fa{خرید} \fa{کردن} (kharid kardan) — to do the shopping, to go shopping}

\label{lesson:FA-C462-kharid-kardan}



\begin{quote}
Before the new one: say the Persian for a sale, then the Persian for to be worth.
\end{quote}

\subsection*{You'll want to know: \fa{خرید} \fa{کردن}}
\textbf{\fa{خرید} \fa{کردن}} (\emph{kharid kardan}) — \textquotedblleft{}to do the shopping, to go shopping\textquotedblright{}.

\begin{quote}
\fa{شنبه} \fa{ها} \fa{مادرم} \fa{در} \fa{بازار} \fa{خرید} \fa{میکند}.
\end{quote}

\begin{quote}
\emph{Shanbe-hâ mâdaram dar bâzâr kharid mikonad.}
\end{quote}

\begin{quote}
On Saturdays my mother does the shopping in the bazaar.
\end{quote}

\emph{kharid} 'purchase' + \emph{kardan}. \emph{Kharidan} is to buy a thing; \emph{kharid kardan} is shopping in general.

\begin{cousinweb}
\emph{kharid kardan} is built from \emph{kharid} 'purchase', the past stem of \emph{kharidan} 'to buy', usually traced to Middle Persian \emph{xrīdan}, and \emph{kardan} 'to do'.
\end{cousinweb}

\begin{grammarlens}[title={What you've built}]
One more word for telling, explaining and describing at B1.
\end{grammarlens}

\subsection*{Your turn}
\begin{itemize}
\raggedright
  \item \emph{Say it:} \emph{kharid kardan}
  \item \emph{Say it:} \emph{kharid kardan}, once more
  \item \emph{Say it:} a sentence of your own with \emph{kharid kardan}
  \item \emph{Recall:} say the Persian for a sale, then the Persian for to be worth, then say \emph{kharid kardan} again
\end{itemize}

\begin{tcolorbox}[breakable,colback=teal!4,colframe=teal!35!black,title={Before you move on}]
What does \fa{خرید} \fa{کردن} mean? (\textquotedblleft{}To do the shopping, to go shopping\textquotedblright{}.) Say it once more.
\end{tcolorbox}
\section[\texorpdfstring{\fa{فروش} \fa{رفتن} (forush raftan)}{فروش رفتن (forush raftan)}]{\fa{فروش} \fa{رفتن} (forush raftan) — to sell (well), to be sold}

\label{lesson:FA-C462-forush-raftan}



\begin{quote}
Before the new one: say the Persian for to be worth, then the Persian for to do the shopping.
\end{quote}

\subsection*{You'll want to know: \fa{فروش} \fa{رفتن}}
\textbf{\fa{فروش} \fa{رفتن}} (\emph{forush raftan}) — \textquotedblleft{}to sell (well), to be sold\textquotedblright{}.

\begin{quote}
\fa{همه} \fa{بلیتها} \fa{در} \fa{یک} \fa{روز} \fa{فروش} \fa{رفت}.
\end{quote}

\begin{quote}
\emph{Hame-ye bilit-hâ dar yek ruz forush raft.}
\end{quote}

\begin{quote}
All the tickets sold in one day.
\end{quote}

\emph{forush} 'sale' + \emph{raftan} 'to go'. \emph{Khub forush miravad} — it sells well.

\begin{cousinweb}
\emph{forush raftan} is built from \emph{forush} 'sale', the present stem of \emph{forukhtan} 'to sell', usually traced to Middle Persian \emph{frōxtan}, and the native verb \emph{raftan} 'to go'.
\end{cousinweb}

\begin{grammarlens}[title={What you've built}]
One more word for telling, explaining and describing at B1.
\end{grammarlens}

\subsection*{Your turn}
\begin{itemize}
\raggedright
  \item \emph{Say it:} \emph{forush raftan}
  \item \emph{Say it:} \emph{forush raftan}, once more
  \item \emph{Say it:} a sentence of your own with \emph{forush raftan}
  \item \emph{Recall:} say the Persian for to be worth, then the Persian for to do the shopping, then say \emph{forush raftan} again
\end{itemize}

\begin{tcolorbox}[breakable,colback=teal!4,colframe=teal!35!black,title={Before you move on}]
What does \fa{فروش} \fa{رفتن} mean? (\textquotedblleft{}To sell (well), to be sold\textquotedblright{}.) Say it once more.
\end{tcolorbox}
\section[\texorpdfstring{\fa{جعبه} (ja'be)}{جعبه (ja'be)}]{\fa{جعبه} (ja'be) — a box, a case}

\label{lesson:FA-C462-jabe}



\begin{quote}
Before the new one: say the Persian for to do the shopping, then the Persian for to sell.
\end{quote}

\subsection*{You'll want to know: \fa{جعبه}}
\textbf{\fa{جعبه}} (\emph{ja'be}) — \textquotedblleft{}a box, a case\textquotedblright{}.

\begin{quote}
\fa{فروشنده} \fa{کفشها} \fa{را} \fa{در} \fa{جعبه} \fa{گذاشت}.
\end{quote}

\begin{quote}
\emph{Forushande kafsh-hâ râ dar ja'be gozâsht.}
\end{quote}

\begin{quote}
The shopkeeper put the shoes in a box.
\end{quote}

\emph{Ja'be-ye shirini} — a box of pastries, a common gift.

\begin{cousinweb}
\emph{ja'be} was borrowed from Arabic \emph{jaʿba} (EtymDB), which in Arabic is a quiver for arrows; in Persian it came to mean any box.
\end{cousinweb}

\begin{grammarlens}[title={What you've built}]
One more word for telling, explaining and describing at B1.
\end{grammarlens}

\subsection*{Your turn}
\begin{itemize}
\raggedright
  \item \emph{Say it:} \emph{ja'be}
  \item \emph{Say it:} \emph{ja'be}, once more
  \item \emph{Say it:} a sentence of your own with \emph{ja'be}
  \item \emph{Recall:} say the Persian for to do the shopping, then the Persian for to sell, then say \emph{ja'be} again
\end{itemize}

\begin{tcolorbox}[breakable,colback=teal!4,colframe=teal!35!black,title={Before you move on}]
What does \fa{جعبه} mean? (\textquotedblleft{}A box, a case\textquotedblright{}.) Say it once more.
\end{tcolorbox}
